% !TEX TS-program = lualatex
Алгоритм <<Альфа-бета-обрезка>> является производительной модификацией алгоритма <<Минимакс>> \cite{alpha-beta-heuristic}. Исходная версия алгоритма требует анализа всего дерева игры для оценки оптимального продолжения, тогда как алгоритм <<Альфа-бета-обрезка>> позволяет двукратно уменьшить сложность (в лучшем случае, если наиболее вероятные сильные ходы рассматриваются первыми) за счёт специальных параметров $\alpha$ и $\beta$. Эти параметры сохраняются в качестве пороговых значений, определяя направление поиска и отсекая критически неоптимальные ветви игрового дерева.

Алгоритм реализуется путём поддержания двух переменных:
\begin{itemize}
    \item $\alpha$ --- наилучшее (максимальное) значение, которое может гарантировать текущий игрок (максимизатор) на текущем пути;
    \item $\beta$ --- наилучшее (минимальное) значение, которое может гарантировать противник (минимизатор) на этом же пути.
\end{itemize}

В процессе обхода дерева, если становится очевидно, что текущий путь не может привести к более выгодному исходу, чем уже известный альтернативный путь, то соответствующая ветвь обрезается. Это называется отсечением по альфа-бета границам. Формально можно сказать, что если в какой-либо вершине выполняется условие $\alpha \geq \beta$, дальнейший анализ поддерева становится избыточным, и обход прекращается для данной ветви. Такой подход позволяет избежать анализа заведомо неэффективных ходов, тем самым значительно сокращая количество вычислений.

Кроме того, на практике альфа-бета-обрезка часто сочетается с эвристиками, упорядочивающими ходы (например, с использованием оценочной функции), что позволяет максимизировать эффект от обрезки и ускорить вычисления.
